\chapter{暴胀、结构与初始条件的丧失}

\section{Omega与平坦性的证据}

一位学生关于大\(\Omega\)的提问把讨论带回到Friedmann-Robertson-Walker方程。对于当前时代，即在辐射主导之后很久，膨胀定律写成一个真空能项与一个物质项，而辐射项被略去，因为它现在非常小：
\begin{equation}
H^2=\left(\frac{\dot a}{a}\right)^2=\frac{8\pi G}{3}\left(\rho_0+\rho_m\right)-\frac{k}{a^2}.
\end{equation}
真空项不稀释。物质项则稀释：
\begin{equation}
\rho_m=\frac{M}{a^3}.
\end{equation}
\(\Omega\)是膨胀方程中由不同组分所承担的比例。

除以\(H^2\)。于是归一化后的各项之和为一。采用正则化的记号，
\begin{equation}
\begin{aligned}
\Omega_0&=\frac{8\pi G\rho_0}{3H^2},\\
\Omega_m&=\frac{8\pi G\rho_m}{3H^2},\\
\Omega_k&=-\frac{k}{a^2H^2},
\end{aligned}
\end{equation}
因此
\begin{equation}
\Omega_0+\Omega_m+\Omega_k=1.
\end{equation}
\editorialnote{这一段中音频与黑板都较粗糙，讲者对真空贡献在几个名称之间来回切换。此处的记号是在除以\(H^2\)之后对论证作正则化。}
如果空间是平坦的，则\(k=0\)，所以
\begin{equation}
\Omega_0+\Omega_m=1.
\end{equation}
如果物质加真空之和大于一，曲率为正。如果小于一，曲率为负。

观测逻辑于是很直接。测量普通物质与暗物质。测量Hubble膨胀。在极大尺度上以几何方式测量曲率。还要注意到宇宙正在加速。真空项的符号本身并不是\(H\)，而是\(\dot a/a\)趋向常数的倾向，这意味着标度因子趋近指数分支，
\begin{equation}
a(t)\propto e^{Ht}.
\end{equation}
这几项测量相互交叉检验，表明宇宙非常接近平坦，误差在百分之几以内。

\begin{classroomqa}
\lecturetimestamp{00:10:50}
\audiencequestion{物质是否包括暗物质？}
\lecturerresponse{是的。就当前的讨论而言，物质项包括暗物质。也可以加上辐射，但今天这一比例小到可以忽略。}
\end{classroomqa}

\begin{classroomqa}
\lecturetimestamp{00:11:11}
\audiencequestion{除了Hubble常数之外，还有别的量度暗能量的方式吗？}
\lecturerresponse{要点不在于\(H\)本身的值，而在于膨胀趋向指数定律的方式。结合物质测量与平坦性的几何测量，就可以推断出真空项。}
\end{classroomqa}

\section{从平坦性到结构}

为什么宇宙如此平坦？其设想是暴胀。平坦意味着在数十亿、实际上是数百亿光年的尺度上非常大、非常光滑且非常均匀。指数膨胀可以把一个区域拉伸到看起来如此。

但好事也可能做过头。一个完全光滑的宇宙会保持完全光滑。引力不会从精确的对称性中制造出成团性。如果现在存在结构，那么结构必定在更早的时候从微小的不均匀性开始。把现在的宇宙往回追溯，在透明时期的不均匀性大约为\(10^5\)分之一：
\begin{equation}
\frac{\delta\rho}{\rho}\sim 10^{-5}.
\end{equation}
宇宙极其光滑，但并非完美光滑。

从那里起，引力做了满屋子空气中的普通力所不做的事。一个小的过密区把物质拉向自身。附近区域被抽空，于是欠密区变得更欠密，过密区变得更过密。引力强化了反差，而不是抹去它。

自然的问题是，为什么这没有更早发生。答案是压强。更早的时候，辐射压很大。利用
\begin{equation}
p=w\rho,
\end{equation}
熟悉的情形是
\begin{equation}
w_{\rm rad}=\frac13,\qquad
w_m=0,\qquad
w_{\rm vac}=-1.
\end{equation}
正压强使介质均匀化。它把扰动推出去并使介质变光滑。物质在极大尺度上被当作无压强的，不提供这样的支撑。有压强存在时，引力被挡住。当压强变得可忽略时，引力获胜。

暗物质在事件的先后顺序中至关重要。最先坍缩的团块是暗物质晕。普通重子物质随后落入这些晕中，并且由于重子可以耗散，会聚集到更靠近中心处。暗物质不以同样的方式耗散，因此它仍然散布在整个晕中。这一收缩的开始比透明时期略早，在时间上大约相差十倍，这就是微波背景没有直接显示出其开端的原因。

一段简短的题外话确定了后来的化学历史。几乎所有比氦重的元素都是恒星演化的产物，而不是早期宇宙宇宙学的产物。早期宇宙基本上只产生氢、氦以及少量几种轻核的痕迹。

在天空中观测到的并不是质量密度本身，而是微波背景中的温度各向异性。密度涨落对它有所贡献。速度涨落也有贡献。过密区域往往看起来略微更冷，因为光子从更深的引力势阱中爬出时会损失额外能量，而气体运动产生的多普勒频移又叠加在其上。完整的数值分析很复杂。稳健的陈述则更简单：到宇宙变得透明时，能量密度的变化约为\(10^5\)分之一。

\begin{classroomqa}
\lecturetimestamp{00:20:29}
\audiencequestion{为什么引力没有让团块更早地增长？}
\lecturerresponse{因为更早的时候宇宙是辐射主导的，而辐射携带了足够的正压来使介质均匀化。只有当辐射与物质相比变得不重要时，引力才开始主导团块的增长。}
\end{classroomqa}

\begin{classroomqa}
\lecturetimestamp{00:24:35}
\audiencequestion{坍缩不需要耗散机制吗？}
\lecturerresponse{普通重子物质若要向中心沉降，确实需要耗散。但暗物质可以在没有那种碰撞耗散的情况下先形成晕。还有一种与膨胀本身相关的宇宙学阻尼，这将在后面变得重要。}
\end{classroomqa}

\begin{classroomqa}
\lecturetimestamp{00:30:34}
\audiencequestion{是否也存在速度涨落，从而使后来的星系能够获得角动量？}
\lecturerresponse{是的。原初图样涉及密度涨落和速度涨落。在天空中看到的是由密度变化和气体运动共同产生的温度图样。}
\end{classroomqa}

\section{暴胀子与平台}

现在引入的机制是一个标量场\(\phi\)。它的能量密度有三部分：
\begin{equation}
\rho_\phi=\frac12\dot\phi^2+\frac12(\nabla\phi)^2+V(\phi).
\end{equation}
时间变化有贡献。空间变化有贡献。标量场还可以携带真正的势能\(V(\phi)\)。这一部分并不奇特；希格斯场就是熟悉的例子。这里重要的是势的形状。

势必须在某个范围内相当平坦，带有轻微向下的斜率，然后在某个临界值处底部塌落。它并不会恰好降到零。它会降到非常小的某个值。在普朗克单位下，当今的真空能约为
\begin{equation}
10^{-123}.
\end{equation}
在平台上时它比这大得多，但仍远低于普朗克尺度。

\begin{figure}[t]
\centering
\includegraphics[width=0.78\textwidth]{lecture_08_figure_03.png}
\caption{具有宽阔平台、陡峭下降和低极小值的定性暴胀子势。}
\lectureframe{00:43:23}
\end{figure}

这个图并不描述空间中的某个位置。它描述的是场的取值。说宇宙从平台高处开始，只是说场值在一个非常大的区域内位于那里。起初\(\phi\)是否随位置有所变化并不重要，因为暴胀会把这些变化极大地拉伸。经过足够多的膨胀之后，任何局部区域看起来都几乎均匀。\editorialnote{黑板画面中的坐标轴标签无法辨认，但口头讨论清楚表明，纵轴是势能\(V(\phi)\)，横轴是场值\(\phi\)。}

只要场停留在势的宽阔、缓慢倾斜的部分上，势能的行为就像真空能。如果\(\phi\)被完全冻结，它就会完全像真空能那样行为。撇开省略的数值因子不谈，讲座中陈述的关系简单地就是
\begin{equation}
H\propto \sqrt{V(\phi)},
\end{equation}
而这给出近乎指数的规律
\begin{equation}
a(t)\propto e^{Ht}.
\end{equation}
现在能够观测到的几乎一切都发生在场越过边缘之后。下降处正是平台上巨大的类真空能转化为辐射、普通物质、暗物质以及今天所见的微小剩余真空能的地方。

\begin{classroomqa}
\lecturetimestamp{00:44:07}
\audiencequestion{\(\phi\)是什么场，爱因斯坦方程仍在使用吗？}
\lecturerresponse{是的。\(\phi\)是时空上的场，同一个弗里德曼方程仍在使用。在这个区域中，只要场变化得足够慢，右侧的能量密度主要由\(V(\phi)\)提供。}
\end{classroomqa}

\begin{classroomqa}
\lecturetimestamp{00:47:04}
\audiencequestion{类真空的暴胀子能量如何转化为物质和辐射？}
\lecturerresponse{当场滚到底部并振荡时，这些振荡就是场的量子。如果这些量子不稳定，它们可以衰变为光子、电子、正电子以及其他熟悉的粒子。}
\end{classroomqa}

\section{慢滚与哈勃摩擦}

一旦场越过边缘，为什么它不会就这样永远晃荡下去？为什么它不会来回摆动，甚至可能再次爬升？这里给出的答案是宇宙学阻尼机制。它不是真正的微观摩擦，但在运动方程中表现得就像摩擦一样。

忽略空间梯度，把场当作均匀的。于是动能项和势能项看起来像一个具有\(T-V\)动力学的单一自由度，但由于它们是能量密度，必须乘以物理体积\(a^3\)。由此得到如下形式的均匀有效拉格朗日量
\begin{equation}
L_{\mathrm{hom}}=a^3(t)\left[\frac12\dot\phi^2-V(\phi)\right].
\end{equation}
\editorialnote{黑板上最上面一行被部分裁掉了。这个均匀拉格朗日量只是被重构到足以重现其下方写出的可靠运动方程的程度。}

黑板上可靠写出的方程是
\begin{equation}
\frac{d}{dt}\left(a^3\dot\phi\right)=-a^3\frac{\partial V}{\partial\phi}.
\end{equation}
展开时间导数得到
\begin{equation}
a^3\ddot\phi+3a^2\dot a\,\dot\phi=-a^3\frac{\partial V}{\partial\phi},
\end{equation}
再除以\(a^3\)，得到
\begin{equation}
\ddot\phi+3H\dot\phi=-\frac{\partial V}{\partial\phi}.
\end{equation}
负号在口头推导中被明确修正：力等于势能导数的负值。

\begin{figure}[t]
\centering
\includegraphics[width=0.78\textwidth]{lecture_08_figure_05.png}
\caption{均匀暴胀子运动方程，其中显示了哈勃摩擦项\(3H\dot\phi\)。}
\lectureframe{00:57:10}
\end{figure}

\(3H\dot\phi\)这一项就是所谓的哈勃摩擦。它完全来自对因子\(a^3\)求导。然而从动力学上看，它就像石块在水中下落时所受的阻力项。这里有一个力项、一个加速度项，以及一个与速度成正比的项。场很快达到终端速度，然后沿着平台缓慢移动。这就是慢滚 regime。当斜率很浅、场移动缓慢时，势能变化也很缓慢，因此膨胀保持接近指数式。

在最小值附近，场发生振荡。这些振荡就是暴胀子量子。如果这些量子不稳定，它们就会衰变，而平台上的能量正是通过这种方式转化为后来宇宙的热物质和辐射内容。

\begin{classroomqa}
\lecturetimestamp{00:50:48}
\audiencequestion{为什么场不会就这样永远振荡下去，又是什么提供了阻尼？}
\lecturerresponse{膨胀本身提供了一个类似摩擦的项。均匀拉格朗日量中的因子\(a^3\)产生了\(3H\dot\phi\)项，而它的行为就像黏性阻力。}
\end{classroomqa}

\section{量子涨落与视界退出}

现在出现了一个表面上的问题。如果暴胀使宇宙变得极其平滑，为什么它没有抹掉星系所需要的那些种子？答案是量子力学。幸存下来的涨落并不是不完美熨平过程的残余。它们是重新产生的。

暂时忘掉膨胀。即使一个场几乎完全均匀，量子力学也不会允许\(\phi\)和\(\dot\phi\)同时被精确固定。基态仍然在扰动。这与普通谐振子具有零点能是同一个道理。在日常情形下，这些涨落很难被看到，因为它们波长短、振荡快，所以观测会对它们取平均。

膨胀改变了情况。在共动坐标中，场的模式大致保持相同的坐标波长，而物理波长则随尺度因子增长：
\begin{equation}
\lambda_{\rm phys}=a(t)\,\Delta x.
\end{equation}
在暴胀期间，物理哈勃尺度近似固定，
\begin{equation}
d_H\sim H^{-1},
\end{equation}
因此相应的共动视界按如下方式收缩
\begin{equation}
d_{H,\mathrm{coord}}\sim \frac{1}{aH}.
\end{equation}

绳子的类比使冻结变得直观。波之所以振荡，是因为绳子的每一小段都受到相邻小段的拉扯。但如果波长变得和视界一样大，相关的相邻部分就超出了因果接触范围。恢复性的拉扯消失了。一旦某个模式比视界更长，它就冻结了。

这种情况一次又一次地发生。较长的波长先冻结。较短的波长后冻结。即使一个模式一开始波长极小，最终也会被拉伸到越过视界。因此，暴胀并不只是抹掉结构。它不断产生冻结的场变化。

这个过程也不会仅仅因为暴胀持续很长时间就停止。旧的冻结模式被拉伸到极大的尺度，但新的模式总是在不断冻结出来。累积起来的图样在宇宙学相关的意义上变得统计上稳定。这就是近似尺度不变谱的起源：同样的视界穿越过程在相继的对数尺度上重复发生，并在相邻波长带之间产生大致相同的振幅。

场滚下平台之后，视界再次增长。曾经处于暂时的暴胀视界之外的模式，后来又会重新进入视界。用通常的说法，它们在暴胀期间离开视界，之后又重新进入视界。到那时它们已被拉伸到宏观尺度，所以当它们再次开始振荡时，是以非常大的结构的形式振荡。

\begin{classroomqa}
\lecturetimestamp{01:12:02}
\audiencequestion{在时空图上，视界看起来是什么样子？}
\lecturerresponse{在指数膨胀期间，到视界的物理距离大致固定在\(c/H\)，但在共动图上，这个固定的物理距离看起来越来越小，因为坐标尺度正在被拉伸。用共动的话说，视界像\(1/(aH)\)那样收缩。}
\end{classroomqa}

\begin{classroomqa}
\lecturetimestamp{01:15:59}
\audiencequestion{当涨落被拉伸时，振幅会发生什么变化，是否存在稳态？}
\lecturerresponse{各个模式在被拉伸时会被稀释，但新模式不断冻结出来。结果是形成一种统计上稳定的图案，近似尺度不变，而不是涨落完全消失。}
\end{classroomqa}

\section{从冻结的场变化到密度变化}

一旦场在不同地点获得了微小的变化，当不同区域以略微不同的时间越过边缘时，这些变化就转化为密度差异。稍微更靠下坡的区域先到达悬崖。如果它的真空能转化为辐射，那辐射会随膨胀而稀释。仍然位于平台上的邻近区域仍以真空能为主，不会稀释。等到第二个区域落下时，第一个区域已经将其一部分能量红移掉了。

这就是为什么滚动越慢，最终的密度反差越大。如果场沿平台移动得非常缓慢，那么场值的微小差异就对应于一个区域与下一个区域到达悬崖之间更大的时间滞后。先落下的区域有更多时间在较晚落下的区域转化其能量之前稀释。坡度越平缓，终端速度越小，最终的密度变化越大。更高的平台也倾向于给出更大的涨落。平台的宽度主要控制暴胀持续多久。

涨落数据只对高度与坡度的某种组合给出严格限制。如果能量尺度足够高，还应当有可见的引力波贡献。这些数值界限只作为粗略的数量级估计给出。粗略地说，如果平台上的能量密度远高于约\(10^{-14}\)（以普朗克单位计），微波背景就已经会显示出过多的引力波结构；如果远低于约\(10^{-16}\)，则看到这种信号的实际机会就很小。\editorialnote{讲者明确将这些数值视为近似值，并说他记不清精确数值。这里保留它们仅作为数量级的陈述。}

单位按通常的自然单位方式解释。质量为\(m\)的粒子的康普顿波长为
\begin{equation}
\lambda_C=\frac{\hbar}{mc}.
\end{equation}
如果\(c=\hbar=1\)，那么质量与长度的倒数具有相同的单位。能量密度是单位体积的能量，所以它具有质量的四次方的单位：
\begin{equation}
[\rho]=(\mathrm{mass})^4.
\end{equation}

在单个点上看，冻结模式的累积看起来像随机行走。空间中不同的点经历不同的随机行走，而这正是把量子噪声转化为空间变化的原因。

\begin{classroomqa}
\lecturetimestamp{01:33:04}
\audiencequestion{\(\phi\)会进行随机行走吗，这会使更早的起点变得无关吗？}
\lecturerresponse{在任意固定点，相继冻结的模式以随机相位相加，所以\(\phi\)确实进行随机行走。正是这种成功也带来了困难：暴胀不断用后来的涨落替换更早的信息，因此暴胀之前的阶段变得难以恢复。}
\end{classroomqa}

\section{暴胀所隐藏的东西}

这个理论在一个非常特定的意义上取得了成功。它把跨越巨大尺度范围的涨落与一个底层机制关联起来。但同一个机制也抹掉了此前存在的东西。更早的不规则性被拉伸掉，而后来产生的量子涨落才是仍然相关的那些。这就是为什么暴胀既吸引人又令人沮丧。它解释了平滑性和涨落谱，但它也抹去了大部分更早的历史。

这导致了对起点的不安。非常高的能量密度并不是问题所在。问题在于，平台能量似乎既不是自然的普朗克尺度，也不是自然的小量。它在普通单位下很高，但在普朗克单位下仍然很小。它看起来并不像是一个基础理论显而易见的第一个停靠点。

在这个简单图像中也不存在字面意义上的\(t=0\)。一旦暴胀持续得过久，关于任何更早阶段的信息几乎全部消失。如果e折数仅仅略高于所需的最小值，也许前一阶段的某些痕迹还能幸存，但一般而言暴胀会抹去过去。

同样的谨慎也出现在对德西特空间的处理中。纯指数律对于晚期暴胀分支是正确的，但精确的德西特解用
\begin{equation}
a(t)\propto \cosh(Ht).
\end{equation}
来描述更好。该解收缩到最小尺寸，然后膨胀。收缩分支在宇宙学中是否应当被字面理解，这一点仍悬而未决。

\begin{classroomqa}
\lecturetimestamp{01:37:31}
\audiencequestion{是否可能没有起点，指数行为无限地延续到过去？}
\lecturerresponse{在这个简单图像中并非如此。精确的德西特空间并不是一个对所有时间都永恒膨胀的指数，它更像是\(\cosh(Ht)\)，有一个收缩分支和一个膨胀分支。那个更早的分支是否应当被字面理解，仍是一个悬而未决的问题。}
\end{classroomqa}

\section{调节、曲率与人类学界限}

在一级近似下，重要的参数是平台的高度、它的斜率和它的宽度\(\Delta\phi\)。斜率重要，因为它决定了终端速度，从而决定了涨落的大小。宽度重要，因为它有助于决定暴胀持续多久。一种常见的偏见是，基本的场偏移不应远大于普朗克质量，并且能量密度不应超过普朗克单位下的量级一。在这样的粗略界限下，足够的暴胀并不能填满整个参数空间。这里引用的估计是百分之几的调节，约为\(2\%\)，以使斜率足够平缓，从而获得大约\(60\)到\(62\)个e折。

\begin{classroomqa}
\lecturetimestamp{01:46:02}
\audiencequestion{暴胀能否用参数空间来理解，它需要多少调节？}
\lecturerresponse{可以。在一级近似下，人们从平台的高度、斜率和宽度来思考。讲座给出了一个粗略估计：获得足够的e折需要百分之几量级的调节，约为\(2\%\)。}
\end{classroomqa}

然后，这种轻微的调节与当前宇宙学常数这一尖锐得多的问题相比较，后者约为\(10^{-123}\)。讨论转向人类学问题，因为这两个问题都可以通过追问结构形成需要什么来重新表述。

在牛顿图像中，正的宇宙学常数的作用就像一个与距离成正比的普遍排斥。如果它大得多，它不必在今天把星系炸开才致命。它只需要大到足以阻止原初\(10^{-5}\)涨落凝聚成星系即可。所给出的例子是，大约比观测值高两个数量级的值就已经能做到这一点。这就是与温伯格论证相关的要点：宇宙学常数不必恰好为零；它只需要小到足以允许结构存在。

负曲率的作用方向相同。在先前讲座的牛顿语言中，\(k<0\)意味着退行速度超过逃逸速度。过大的向外推力也会阻止微小的原初超密度最终变成星系。暴胀稀释了那种曲率。所引用的下界约为\(59.5\)个e折。然后，大约\(62\)的值被当作一个安全高于阈值的代表性暴胀历史。

这改变了概率问题。如果人们追问给出足够暴胀的参数空间的原始体积，答案可能只有百分之几。如果人们以身处一个星系形成、观察者存在从而能提出这个问题的宇宙为条件，那么条件概率可以变得大得多。讨论就在这里结束：不是以一个最终解决方案结束，而是以微调问题、对它们的人类学重新解释，以及认识到伟大争论仍然在此处这一事实结束。
